\documentclass[11pt]{article}
\usepackage{amsmath,amssymb,amsfonts,amsthm}
\usepackage{geometry}
\usepackage{hyperref}
\usepackage{bm}
\usepackage{graphicx}
\usepackage{enumitem}
\geometry{margin=1in}

\title{Endogenous Spectral Self-Regulation in Tensor-Extended DAEs\\
and Multiplicity-Driven Feedback:\\
A Defensive Publication}
\author{Ryan O. Van Gelder\\Citizen Gardens -- The Foundation of Multiplicity}
\date{May 5, 2026}

\begin{document}
\maketitle

\begin{abstract}
This defensive publication documents a framework that extends differential-algebraic equations (DAEs) with tensor representations, dynamic eigenvalue multiplicities, and Lyapunov-based spectral feedback. It formalizes how a spectral observable---constructed from the eigenvalue gap of a contracted tensor block---can act as an endogenous state that drives its own stabilization through passivity-preserving feedback and multiplicity-shaped dissipation. We present (1) a generalized tensor-parameterized DAE structure, (2) multiplicity operators and feedback laws, (3) a minimal two-oscillator benchmark with analytically tractable Lyapunov functions, and (4) explicit numerical implementations and diagnostics. The document is intended as prior art and a comprehensive technical record of the concepts and constructions developed in this thread.
\end{abstract}
\newpage
\tableofcontents

\section{Executive Summary}

\subsection{Conceptual Contributions}

This work discloses a family of dynamical systems in which \emph{eigenvalue multiplicity} is promoted from a static algebraic invariant to a \emph{dynamic control variable} that participates in closed-loop feedback.

The key contributions are:

\begin{enumerate}[label=(\alph*)]
  \item A tensor-extended DAE framework
  \[
    \frac{d}{dt}\left(T(t,z)\cdot z\right)
    = \big(J(t,z) - R(t,z)\big)\,Q(t,z)\,z + B(t)u,
  \]
  where $T(t,z)$ is an $n\times n\times k$ tensor encoding higher-order couplings rather than a static mass matrix.

  \item A multiplicity observable $M$ built from the spectrum of a contracted block $M_{\text{mass}}(\kappa)$, which serves as an endogenous state variable that influences dissipation and feedback.

  \item Lyapunov-based feedback laws in which a spectral quantity (typically the eigenvalue gap) defines a Lyapunov function $H$, and the feedback on coupling parameters (e.g.\ $\kappa$) is taken as the negative gradient $-\nabla H$.

  \item A minimal two-oscillator mechanical benchmark where:
  \begin{itemize}
    \item the contracted tensor block is invertible and well-conditioned along the trajectory,
    \item the Lyapunov derivative $\dot H$ is nonpositive for all times (passivity),
    \item the dissipation coefficient $c(t)$ is driven by the rate of change of multiplicity $|\dot M(t)|$ and smoothly returns to a baseline as the system settles.
  \end{itemize}
\end{enumerate}

Together, these constructions realize what we call \emph{endogenous spectral self-regulation}: the system senses its own spectral configuration and adjusts dissipative dynamics accordingly, with formal Lyapunov guarantees.

\subsection{Defensive Publication Scope}

This document should be read as a defensive publication covering, in particular:

\begin{itemize}
  \item Use of a tensor-valued ``mass'' object $T(t,z)$ in DAEs, inverted in practice via monitored pseudoinverses on contracted slices, with rank and condition-number diagnostics integrated into the dynamics.
  \item Construction of smooth multiplicity observables $M(\Delta)$ as functions of eigenvalue gaps, serving as dynamic states that drive feedback.
  \item Lyapunov functions defined on spectral gaps $\Delta$ (or multiplicity $M$) with feedback laws of the form $\dot\kappa = -\nabla_\kappa H(\Delta(\kappa))$, guaranteeing $\dot H \le 0$.
  \item Coupling dissipation $R$ or scalar damping $c(t)$ to the rate of change $|\dot M(t)|$ to create multiplicity-shaped dissipation.
  \item Concrete implementation details and diagnostics for a two-oscillator benchmark (state structure, ODEs, numerical scheme, rank and conditioning measures).
\end{itemize}

\section{Generalized Tensor-Extended DAE Framework}

\subsection{Baseline DAE and Dissipative Port-Hamiltonian Form}

Classical DAEs in dissipative port-Hamiltonian form can be written as
\begin{equation}
  \frac{d}{dt}(E z) = (J - R) Q z + B u,
\end{equation}
with:
\begin{itemize}
  \item $E \in \mathbb{R}^{n\times n}$: (possibly singular) mass matrix.
  \item $J\in \mathbb{R}^{n\times n}$: skew-symmetric energy-exchange matrix.
  \item $R\in \mathbb{R}^{n\times n}$: symmetric positive semidefinite dissipation matrix.
  \item $Q$: energy-storage matrix (often related to $\nabla^2 H$).
  \item $B$: input map, $u$: external input.
\end{itemize}

This structure underlies a large class of energy-based models in control, mechanics, and networked systems.

\subsection{Tensor-Parameterized Dynamics}

We generalize $E$ to a tensor-valued map:
\begin{equation}
  \frac{d}{dt}\big(T(t,z)\cdot z\big)
  = \big(J(t,z) - R(t,z)\big)\,Q(t,z)\,z + B(t)u,
\end{equation}
where:
\begin{itemize}
  \item $T(t,z)\in\mathbb{R}^{n\times n\times k}$ encodes higher-order couplings (multi-scale, multi-body).
  \item The contraction $T(t,z)\cdot z$ is defined via a fixed set of weights $w_\ell(t)$:
  \[
    [T(t,z)\cdot z]_i
    = \sum_{j=1}^n \sum_{\ell=1}^k T_{ij\ell}(t,z)\,z_j\,w_\ell(t),
  \]
  producing an effective $n\times n$ block per time $t$.
  \item The ``inverse'' $T^{-1}(t,z)$ in numerical solvers is implemented as the Moore–Penrose pseudoinverse of this contracted block, with continuous monitoring of:
  \begin{align*}
    \text{rank}\,T_{\text{eff}}(t,z), \quad
    \kappa(T_{\text{eff}}(t,z)) = \|T_{\text{eff}}\|\,\|T_{\text{eff}}^{+}\|,
  \end{align*}
  to avoid singular-mass failure modes.
\end{itemize}

This realizes a tensor-extended DAE in which the ``mass'' object is not static but driven by time, state, and higher-order structure.

\section{Spectral Multiplicity and Feedback}

\subsection{Multiplicity Operator from Eigenvalue Gap}

Let $M_{\text{mass}}(\kappa) \in \mathbb{R}^{2\times 2}$ denote a contracted mass block depending on a scalar coupling parameter $\kappa$:
\begin{equation}
  M_{\text{mass}}(\kappa) =
  \begin{bmatrix}
    m_1 & \kappa\\
    \kappa & m_2
  \end{bmatrix}.
\end{equation}
Its eigenvalues are
\begin{equation}
  \lambda_{\pm}(\kappa) = \frac{(m_1+m_2)\pm \Delta(\kappa)}{2},
\end{equation}
where the gap
\begin{equation}
  \Delta(\kappa) = |\lambda_+ - \lambda_-|
  = \sqrt{(m_1-m_2)^2 + 4\kappa^2}.
\end{equation}

We construct a smooth multiplicity observable $M(\Delta)$ as a function of the gap. Two canonical forms are:

\begin{enumerate}[label=(\roman*)]
  \item ``Plateau'' multiplicity (used in the initial spectral benchmark):
  \begin{equation}
    M(\Delta) = 2 - \exp\left(-\frac{\Delta^2}{2\sigma^2}\right),
  \end{equation}
  which tends to $1$ at exact degeneracy ($\Delta=0$) and to $2$ as $\Delta\to\infty$.
  \item Degeneracy-seeking multiplicity (preferred for visible transients):
  \begin{equation}
    M(\Delta) = 1 + \exp\left(-\frac{\Delta^2}{2\sigma^2}\right),
  \end{equation}
  where $M\to 2$ as $\Delta\to 0$ (double eigenvalue) and $M\to 1$ as $\Delta$ becomes large.
\end{enumerate}

In both cases, $M(\Delta)$ is smooth, bounded, and differentiable with respect to $\kappa$ via the chain rule.

\subsection{Derivatives and Chain Rule}

For either exponential form,
\begin{equation}
  \frac{\partial M}{\partial \Delta}
  = \pm \frac{\Delta}{\sigma^2}
     \exp\left(-\frac{\Delta^2}{2\sigma^2}\right),
\end{equation}
with a sign determined by the chosen offset (the algebraic structure is the same), and
\begin{equation}
  \frac{\partial \Delta}{\partial \kappa}
  = \frac{4\kappa}{\Delta}
  \quad \text{for } \Delta > 0,
\end{equation}
so that
\begin{equation}
  \frac{\partial M}{\partial \kappa}
  = \frac{\partial M}{\partial \Delta}
    \frac{\partial \Delta}{\partial \kappa}
  = \pm \frac{4\kappa}{\sigma^2}
      \exp\left(-\frac{\Delta^2}{2\sigma^2}\right).
\end{equation}

If we choose a feedback law $\dot\kappa = g(\kappa)$, then
\begin{equation}
  \dot M
  = \frac{\partial M}{\partial \kappa}\,\dot\kappa
  = \pm \frac{4\kappa}{\sigma^2}
      \exp\left(-\frac{\Delta^2}{2\sigma^2}\right)\,g(\kappa).
\end{equation}
This provides a closed analytic expression linking the multiplicity dynamics to coupling dynamics.

\section{Lyapunov-Based Spectral Feedback}

\subsection{Gap-Based Lyapunov Function}

A natural Lyapunov function for spectral regulation is
\begin{equation}
  H(\Delta) := \frac{1}{2}\Delta^2.
\end{equation}
In the symmetric-mass case $m_1 = m_2$, we have $\Delta(\kappa) = 2|\kappa|$, so
\begin{equation}
  H(\kappa) = \frac{1}{2}(2|\kappa|)^2 = 2\kappa^2.
\end{equation}

A gradient-descent feedback law in $\kappa$ is
\begin{equation}
  \dot\kappa = -\beta\,\frac{\partial H}{\partial \kappa}
  = -\beta \cdot 4\kappa = -4\beta\kappa,
\end{equation}
leading to an exponential relaxation:
\begin{equation}
  \kappa(t) = \kappa_0\,e^{-4\beta t}.
\end{equation}
The Lyapunov derivative is
\begin{equation}
  \dot H = \frac{\partial H}{\partial \kappa}\,\dot\kappa
  = 4\kappa \cdot (-4\beta\kappa)
  = -16\beta\kappa^2 \le 0.
\end{equation}
Thus $H$ is strictly decreasing except at $\kappa=0$.

\subsection{Relaxation Timescale}

For $H = 2\kappa^2$, the condition $H(t) \le \alpha H(0)$ is equivalent to
\begin{equation}
  \kappa^2(t) \le \alpha \kappa_0^2,
  \quad \Rightarrow\quad
  e^{-8\beta t} \le \alpha,
\end{equation}
yielding
\begin{equation}
  t_\alpha = \frac{\ln(1/\alpha)}{8\beta}.
\end{equation}
For $\alpha = 0.05$ (a 95\% reduction), this gives
\begin{equation}
  t_{95}
  = \frac{\ln 20}{8\beta}.
\end{equation}
This analytical timescale can be compared against numerical estimates based on interpolating $H(t)$.

\section{Multiplicity-Shaped Dissipation}

\subsection{Dissipation as a Function of \texorpdfstring{$\dot M$}{dot M}}

We define a scalar dissipation coefficient $c(t)$ as
\begin{equation}
  c(t) = c_0 + \gamma |\dot M(t)|,
\end{equation}
where $c_0$ is a baseline damping and $\gamma>0$ scales the sensitivity to multiplicity change. Using the earlier expression for $\dot M$, we get
\begin{equation}
  c(t) = c_0 + \gamma \left|
    \frac{\partial M}{\partial \kappa} \dot\kappa
  \right|
  = c_0 + \gamma
    \left|
      \pm \frac{4\kappa}{\sigma^2}
          \exp\left(-\frac{\Delta^2}{2\sigma^2}\right)\cdot (-4\beta\kappa)
    \right|.
\end{equation}
This is largest when both $|\kappa|$ and the exponential factor are substantial; it decays as $\kappa\to 0$ or as the system moves into a regime where the Gaussian factor suppresses changes in $M$.

\subsection{Qualitative Behavior}

The intended qualitative behavior is:

\begin{itemize}
  \item For large spectral gap (far from degeneracy), $M$ is close to 1, $\dot M$ is small, and $c(t)\approx c_0$.
  \item As $\kappa(t)$ and hence $\Delta(t)$ evolve into the sensitive region of the Gaussian, $M(t)$ changes appreciably, $|\dot M(t)|$ rises, and $c(t)$ experiences a transient increase (a ``dissipative bump'').
  \item As the system settles into a regime where $\Delta$ stabilizes (or reaches degeneracy in the symmetric case), $\dot M\to 0$ and $c(t)\to c_0$ smoothly.
\end{itemize}

This realizes multiplicity-shaped dissipation that is maximally active during significant spectral reconfiguration and quiescent in static spectral regimes.

\section{Two-Oscillator Benchmark: Dynamics and Implementation}

\subsection{Mechanical Setup}

Consider two linearly coupled oscillators with positions $q_1, q_2$ and momenta $p_1, p_2$, masses $m_1, m_2$, and frequencies $\omega_1, \omega_2$. Let $z = [q_1, p_1, q_2, p_2]^\top$.

The contracted mass block is
\begin{equation}
  M_{\text{mass}}(\kappa) =
  \begin{bmatrix}
    m_1 & \kappa\\
    \kappa & m_2
  \end{bmatrix},
\end{equation}
with pseudoinverse $M_{\text{mass}}^{+}$ and state velocities
\begin{equation}
  \begin{bmatrix}
    \dot q_1\\ \dot q_2
  \end{bmatrix}
  = M_{\text{mass}}^{+}(\kappa)
    \begin{bmatrix}
      p_1\\ p_2
    \end{bmatrix}.
\end{equation}

\subsection{Equations of Motion}

The equations of motion in the benchmark are:
\begin{align}
  \dot q_1 &= v_1,\\
  \dot q_2 &= v_2,
\end{align}
where $(v_1,v_2)^\top = M_{\text{mass}}^{+}(\kappa)\,(p_1,p_2)^\top$, and
\begin{align}
  \dot p_1 &= -\big(\omega_1^2 q_1 + \kappa (q_1 - q_2)\big)
             - c(t)\,(v_1 - v_2),\\
  \dot p_2 &= -\big(\omega_2^2 q_2 + \kappa (q_2 - q_1)\big)
             - c(t)\,(v_2 - v_1),
\end{align}
with $c(t)$ as above.

\subsection{Coupling Dynamics and Spectral Feedback}

The coupling parameter evolves via
\begin{equation}
  \dot\kappa = -4\beta\kappa,
\end{equation}
yielding $\kappa(t) = \kappa_0 e^{-4\beta t}$, and the gap-based Lyapunov function
\begin{equation}
  H(t) = \frac{1}{2}\Delta^2(t)
\end{equation}
satisfies $\dot H \le 0$.

In the general asymmetric-mass case ($m_1\neq m_2$), the gap has a nonzero floor, so $H$ asymptotically approaches a positive constant; in the symmetric case ($m_1 = m_2$), the gap can collapse to zero and $H\to0$.

\subsection{Rank and Conditioning Diagnostics}

At each time step, we compute:
\begin{align}
  r(t) &= \text{rank}\,M_{\text{mass}}(\kappa(t)),\\
  \kappa_{\text{cond}}(t)
      &= \|M_{\text{mass}}(\kappa(t))\|\,
         \|M_{\text{mass}}^{+}(\kappa(t))\|,
\end{align}
to ensure that the pseudoinverse remains numerically stable. In the numerical experiments, $r(t)$ remained equal to 2 and $\kappa_{\text{cond}}(t)$ stayed below 2, demonstrating good conditioning.

\section{Representative Code Snippet}

The following Python snippet captures the core of the benchmark: mass contraction, Lyapunov-based coupling dynamics, multiplicity observable, dissipation update, and integration of the mechanical state.

\subsection*{Python Implementation (Illustrative)}

\begin{verbatim}
import numpy as np
from scipy.integrate import solve_ivp

# Parameters
m1, m2 = 1.0, 1.0        # symmetric mass case
omega1, omega2 = 1.0, 1.15
beta = 0.8
gamma = 0.35
c0 = 0.12
sigma = 0.5              # multiplicity sensitivity scale

def mass_block(kappa):
    return np.array([[m1, kappa],
                     [kappa, m2]], dtype=float)

def gap(kappa):
    return np.sqrt((m1 - m2)**2 + 4.0*kappa**2)

def multiplicity(delta, sigma):
    # Degeneracy-seeking form: M -> 1 (large gap), M -> 2 (small gap)
    return 1.0 + np.exp(-delta**2/(2.0*sigma**2))

def dM_dt(kappa, delta, beta, sigma):
    # Chain rule: dM/dt = (dM/dkappa) * (dkappa/dt)
    # dDelta/dkappa = 4*kappa / delta (for delta > 0)
    if delta <= 1e-12:
        return 0.0
    dM_ddelta = (delta / sigma**2) * np.exp(-delta**2/(2.0*sigma**2))
    ddelta_dk = 4.0 * kappa / delta
    dM_dk = dM_ddelta * ddelta_dk
    dk_dt = -4.0 * beta * kappa
    return dM_dk * dk_dt

def rhs_state(t, y, kappa, c):
    q1, p1, q2, p2 = y
    Minv = np.linalg.pinv(mass_block(kappa))
    v1, v2 = Minv @ np.array([p1, p2])
    dp1 = -(omega1**2 * q1 + kappa * (q1 - q2)) - c * (v1 - v2)
    dp2 = -(omega2**2 * q2 + kappa * (q2 - q1)) - c * (v2 - v1)
    return np.array([v1, dp1, v2, dp2])

# Initial conditions
y0 = np.array([1.0, 0.0, -0.4, 0.0], dtype=float)
kappa0 = 1.5
T_end = 8.0
outer_dt = 0.05

t_cur = 0.0
y = y0.copy()
kappa = kappa0

trajectory = []

while t_cur < T_end - 1e-12:
    delta = gap(kappa)
    M_val = multiplicity(delta, sigma)
    dM_val = dM_dt(kappa, delta, beta, sigma)
    c = c0 + gamma * abs(dM_val)

    # Integrate fast mechanical dynamics over [t_cur, t_cur + outer_dt]
    sol = solve_ivp(lambda t, yy: rhs_state(t, yy, kappa, c),
                    (t_cur, min(t_cur + outer_dt, T_end)),
                    y, method='RK45',
                    max_step=0.005,
                    rtol=1e-9, atol=1e-11)
    for tt, yy in zip(sol.t, sol.y.T):
        trajectory.append({
            't': tt,
            'q1': yy[0], 'p1': yy[1],
            'q2': yy[2], 'p2': yy[3],
            'kappa': kappa,
            'delta': delta,
            'M': M_val,
            'c': c
        })

    y = sol.y[:, -1]
    t_cur = sol.t[-1]

    # Update coupling parameter with Lyapunov-based feedback
    kappa = kappa - outer_dt * 4.0 * beta * kappa
\end{verbatim}

This snippet illustrates the key elements:

\begin{itemize}
  \item Tensor/contraction as a mass block with pseudoinverse.
  \item Gap-based multiplicity and its derivative.
  \item Dissipation modulated by $|\dot M|$.
  \item Simple outer-loop update for $\kappa$ consistent with $\dot\kappa = -4\beta\kappa$.
\end{itemize}

\section{Prior Art Context and Novelty}

\subsection{Relation to Classical Lyapunov and Spectral Methods}

Classical Lyapunov theory treats $V(x)$ as a state-space functional whose decrease certifies stability; here, the Lyapunov function is defined not on the full state $z$, but on a \emph{spectral feature} (the gap $\Delta$) of a tensor-contracted mass block.

While Lyapunov-based controllers and spectral analysis are individually well-established, the following aspects appear novel:

\begin{itemize}
  \item Using a \emph{tensor-parameterized mass} with an explicit pseudoinverse-based implementation as the carrier of spectral structure.
  \item Treating \emph{eigenvalue multiplicity} (or smooth proxies thereof) as a dynamic state that directly informs feedback and dissipation.
  \item Encoding multiplicity dynamics in a Lyapunov function defined on the gap, and using its gradient with respect to coupling parameters as the feedback law in a DAE context.
\end{itemize}

\subsection{Defensive Publication Claims}

To make the defensive nature explicit, we highlight:

\begin{enumerate}[label=(\arabic*)]
  \item A DAE framework with state-dependent, tensor-valued mass $T(t,z)$, inverted via monitored pseudoinverses, and integrated into port-Hamiltonian-style dynamics.
  \item Spectral multiplicity observables $M$ constructed from eigenvalue gaps of contracted tensor slices, treated as endogenous, differentiable state variables.
  \item Feedback laws in which coupling parameters (e.g.\ $\kappa$) evolve according to $-\nabla_\kappa H(\Delta(\kappa))$, where $H$ is a Lyapunov function defined on the eigenvalue gap.
  \item Dissipation terms $R(t)$ or scalar damping coefficients $c(t)$ defined as explicit functions of $|\dot M(t)|$, so that multiplicity dynamics directly shape dissipative behavior.
  \item A concrete two-oscillator mechanical implementation with:
  \begin{itemize}
    \item Numerically stable pseudoinverses (rank=2, bounded condition number).
    \item Nonpositive Lyapunov derivative throughout the trajectory.
    \item Dissipation returning smoothly to baseline as the coupling parameter relaxes.
  \end{itemize}
\end{enumerate}

\section{Extensions and Future Work}

The constructions presented here are readily extensible:

\begin{itemize}
  \item Higher-dimensional mechanical or network systems with larger contracted blocks and higher-order tensors $T(t,z)$.
  \item Quantum-inspired models where $M$ is built from gaps in a Hamiltonian spectrum and feeds back into dissipative Lindbladian terms.
  \item Control architectures where multiplicity observables act as safety monitors to avoid dangerous bifurcations (gap stabilization rather than degeneracy seeking).
\end{itemize}

These directions are left as future work but are naturally supported by the general framework and benchmark documented in this publication.

\section*{Mathematical Appendix: Lyapunov Proofs and Operator Bounds}

\subsection*{A. Spectral Structure of the Contracted Mass Block}

We recall the contracted $2\times 2$ mass block
\[
  M_{\text{mass}}(\kappa)
  = \begin{bmatrix}
      m_1 & \kappa\\
      \kappa & m_2
    \end{bmatrix},
\]
with real parameters $m_1, m_2 > 0$ and scalar coupling $\kappa \in \mathbb{R}$.

\subsubsection*{A.1 Eigenvalues and Gap}

The characteristic polynomial is
\[
  p(\lambda)
  = \det\!\big(M_{\text{mass}}(\kappa) - \lambda I\big)
  = (\lambda - m_1)(\lambda - m_2) - \kappa^2.
\]
Solving $p(\lambda) = 0$ gives the eigenvalues
\begin{equation}
  \lambda_\pm(\kappa)
  = \frac{1}{2}(m_1 + m_2)
    \pm \frac{1}{2}\sqrt{(m_1 - m_2)^2 + 4\kappa^2}.
\end{equation}
We define the eigenvalue gap
\begin{equation}
  \Delta(\kappa)
  := |\lambda_+ - \lambda_-|
  = \sqrt{(m_1 - m_2)^2 + 4\kappa^2}.
\end{equation}

\paragraph{Symmetric Mass Case.} If $m_1 = m_2 = m$, then
\[
  \Delta(\kappa) = \sqrt{0 + 4\kappa^2} = 2|\kappa|,
\]
and the eigenvalues simplify to
\[
  \lambda_\pm(\kappa) = m \pm |\kappa|.
\]
In particular, degeneracy ($\lambda_+ = \lambda_-$) occurs if and only if $\kappa = 0$.

\subsubsection*{A.2 Rank and Determinant}

The determinant of $M_{\text{mass}}(\kappa)$ is
\begin{equation}
  \det M_{\text{mass}}(\kappa)
  = m_1 m_2 - \kappa^2.
\end{equation}
Thus:
\begin{itemize}
  \item $M_{\text{mass}}(\kappa)$ is invertible if and only if $m_1 m_2 - \kappa^2 \neq 0$.
  \item Rank deficiency occurs iff $\kappa^2 = m_1 m_2$, in which case one eigenvalue vanishes and the other is $m_1 + m_2$.
\end{itemize}

In the numerical experiments, parameters are chosen so that $\kappa^2 \ll m_1 m_2$ along the trajectory, ensuring
\[
  \det M_{\text{mass}}(\kappa) > 0, \quad \text{rank}\,M_{\text{mass}}(\kappa) = 2.
\]

\subsubsection*{A.3 Norm and Condition Number Bounds}

Let $\|\cdot\|_2$ denote the spectral (operator) norm. Since $M_{\text{mass}}(\kappa)$ is real symmetric, its eigenvalues are real and
\[
  \|M_{\text{mass}}(\kappa)\|_2 = \max\{|\lambda_+(\kappa)|, |\lambda_-(\kappa)|\}.
\]
If $M_{\text{mass}}(\kappa)$ is invertible, then
\[
  \|M_{\text{mass}}(\kappa)^{-1}\|_2
  = \frac{1}{\min\{|\lambda_+(\kappa)|, |\lambda_-(\kappa)|\}}.
\]
The $2$-norm condition number is
\[
  \kappa_2\big(M_{\text{mass}}(\kappa)\big)
  := \|M_{\text{mass}}(\kappa)\|_2
     \|M_{\text{mass}}(\kappa)^{-1}\|_2
  = \frac{\max\{|\lambda_+|,|\lambda_-|\}}
         {\min\{|\lambda_+|,|\lambda_-|\}}.
\]

\paragraph{Symmetric Mass Case.} For $m_1=m_2=m>0$ and $|\kappa|<m$, the eigenvalues are $m\pm|\kappa|$ and hence
\[
  \|M_{\text{mass}}(\kappa)\|_2 = m + |\kappa|,
  \quad
  \|M_{\text{mass}}(\kappa)^{-1}\|_2 = \frac{1}{m - |\kappa|},
\]
giving
\begin{equation}
  \kappa_2\big(M_{\text{mass}}(\kappa)\big)
  = \frac{m + |\kappa|}{m - |\kappa|}.
\end{equation}
In particular, if $|\kappa|\le \alpha m$ for some $0<\alpha<1$, then
\begin{equation}
  \kappa_2\big(M_{\text{mass}}(\kappa)\big)
  \le \frac{1+\alpha}{1-\alpha},
\end{equation}
which is uniformly bounded. This bound underlies the observation that the pseudoinverse remains well-conditioned in the benchmark regime.

\subsection*{B. Lyapunov-Based Gap Regulation}

\subsubsection*{B.1 Lyapunov Function on the Gap}

We define a scalar Lyapunov function
\begin{equation}
  H(\Delta) := \frac{1}{2}\Delta^2.
\end{equation}
By construction, $H(\Delta)\ge 0$ and $H(\Delta)=0$ if and only if $\Delta=0$, i.e.\ at exact degeneracy.

\paragraph{Symmetric Mass Case.} For $m_1=m_2=m$, we have $\Delta(\kappa)=2|\kappa|$, hence
\begin{equation}
  H(\kappa) = \frac{1}{2}(2|\kappa|)^2 = 2\kappa^2.
\end{equation}
We take $H(\kappa) = 2\kappa^2$ as a Lyapunov candidate in the $\kappa$-coordinate.

\subsubsection*{B.2 Feedback Law and Lyapunov Derivative}

Define the feedback law
\begin{equation}
  \dot\kappa = -4\beta\kappa, \quad \beta>0,
\end{equation}
which is equivalent to gradient descent on $H$:
\begin{align}
  \frac{\partial H}{\partial \kappa}
  &= 4\kappa,\\
  \dot\kappa
  &= -\beta\frac{\partial H}{\partial \kappa}
   = -4\beta\kappa.
\end{align}

Then
\begin{align}
  \dot H(\kappa)
  &= \frac{\partial H}{\partial \kappa}\,\dot\kappa
   = 4\kappa \cdot (-4\beta\kappa)
   = -16\beta\kappa^2.
\end{align}
This yields the following:

\paragraph{Proposition B.1 (Lyapunov Decrease).} For the symmetric mass case and feedback law $\dot\kappa=-4\beta\kappa$, the Lyapunov function $H(\kappa)=2\kappa^2$ satisfies
\[
  \dot H(\kappa(t)) \le 0,\quad \forall t\ge 0,
\]
with equality if and only if $\kappa(t)=0$.

\emph{Proof.} Immediate from $\dot H = -16\beta\kappa^2$ and $\beta>0$. \qed

\subsubsection*{B.3 Closed-Form Solution and Relaxation Time}

Solving $\dot\kappa = -4\beta\kappa$ with initial condition $\kappa(0)=\kappa_0$ yields
\begin{equation}
  \kappa(t) = \kappa_0\,e^{-4\beta t}.
\end{equation}
Then
\[
  H(t) = 2\kappa(t)^2 = 2\kappa_0^2 e^{-8\beta t}.
\]
For any $0<\alpha<1$, the time $t_\alpha$ at which $H(t)\le \alpha H(0)$ is given by
\[
  2\kappa_0^2 e^{-8\beta t_\alpha} \le \alpha \cdot 2\kappa_0^2
  \quad\Rightarrow\quad
  e^{-8\beta t_\alpha} \le \alpha
\]
and thus
\begin{equation}
  t_\alpha \ge \frac{\ln(1/\alpha)}{8\beta}.
\end{equation}
For the 95\% reduction level $\alpha = 0.05$, we obtain
\begin{equation}
  t_{95}
  = \frac{\ln 20}{8\beta}.
\end{equation}

\subsection*{C. Multiplicity Observable and Its Derivative}

\subsubsection*{C.1 Smooth Multiplicity from Gap}

Define a smooth multiplicity observable $M:\mathbb{R}_{\ge 0}\to(1,2]$ by
\begin{equation}
  M(\Delta) = 1 + \exp\left(-\frac{\Delta^2}{2\sigma^2}\right),
  \quad \sigma>0.
\end{equation}
Then:
\[
  \lim_{\Delta\to\infty} M(\Delta) = 1,
  \quad
  M(0) = 2,
\]
so multiplicity moves from $1$ (well-separated eigenvalues) toward $2$ (degeneracy) as the gap closes.

\subsubsection*{C.2 Derivative with Respect to Gap and Coupling}

Computing the derivative with respect to $\Delta$,
\begin{align}
  \frac{\partial M}{\partial \Delta}(\Delta)
  &= \exp\left(-\frac{\Delta^2}{2\sigma^2}\right)
     \left(-\frac{\Delta}{\sigma^2}\right)
     \cdot (-1)\\
  &= \frac{\Delta}{\sigma^2}
     \exp\left(-\frac{\Delta^2}{2\sigma^2}\right).
\end{align}
For $\Delta>0$, the derivative of $\Delta(\kappa)$ with respect to $\kappa$ is
\begin{align}
  \frac{\partial \Delta}{\partial \kappa}
  &= \frac{1}{2}\left[(m_1 - m_2)^2 + 4\kappa^2\right]^{-1/2} \cdot 8\kappa
   = \frac{4\kappa}{\Delta}.
\end{align}
Thus,
\begin{align}
  \frac{\partial M}{\partial \kappa}
  &= \frac{\partial M}{\partial \Delta}
     \frac{\partial \Delta}{\partial \kappa}\\
  &= \left(
       \frac{\Delta}{\sigma^2}
       \exp\left(-\frac{\Delta^2}{2\sigma^2}\right)
     \right)
     \left(
       \frac{4\kappa}{\Delta}
     \right)\\
  &= \frac{4\kappa}{\sigma^2}
     \exp\left(-\frac{\Delta^2}{2\sigma^2}\right).
\end{align}

\subsubsection*{C.3 Time Derivative Under Spectral Feedback}

Given the feedback law $\dot\kappa = -4\beta\kappa$, we obtain
\begin{align}
  \dot M
  &= \frac{\partial M}{\partial \kappa}\,\dot\kappa\\
  &= \frac{4\kappa}{\sigma^2}
     \exp\left(-\frac{\Delta^2}{2\sigma^2}\right)
     (-4\beta\kappa)\\
  &= -\frac{16\beta\kappa^2}{\sigma^2}
      \exp\left(-\frac{\Delta^2}{2\sigma^2}\right).
\end{align}
Hence:

\paragraph{Proposition C.1 (Sign and Smoothness of \texorpdfstring{$\dot M$}{dot M}).}
For $\sigma>0$, $\beta>0$ and any $\kappa\in\mathbb{R}$:
\begin{enumerate}
  \item $\dot M(\kappa) \le 0$ with equality if and only if $\kappa=0$ (for the degeneracy-seeking form).
  \item $\dot M(\kappa)$ is smooth in $\kappa$ and tends to $0$ as either $\kappa\to 0$ or $|\kappa|\to\infty$.
\end{enumerate}

\emph{Proof.} The factor $-16\beta\kappa^2/\sigma^2$ is nonpositive, with equality only at $\kappa=0$. The exponential factor is positive and decays to $0$ as $|\Delta|\to\infty$; continuity in $\kappa$ follows from continuity of all components. \qed

\subsection*{D. Multiplicity-Shaped Dissipation}

\subsubsection*{D.1 Definition and Bounds}

Define
\begin{equation}
  c(t) = c_0 + \gamma |\dot M(t)|
       = c_0 + \gamma \left|
         -\frac{16\beta\kappa^2(t)}{\sigma^2}
         \exp\left(-\frac{\Delta^2(t)}{2\sigma^2}\right)
       \right|,
\end{equation}
with $c_0>0$, $\gamma>0$.

Then:
\[
  c(t) \ge c_0,
\]
and
\begin{equation}
  c(t) \le c_0 + \gamma \frac{16\beta}{\sigma^2}\kappa^2(t),
\end{equation}
since $0<\exp(-\Delta^2/(2\sigma^2))\le 1$.

\paragraph{Symmetric Mass Case.} With $m_1=m_2$, we have $\Delta = 2|\kappa|$, and
\begin{equation}
  c(t) = c_0 + \gamma \frac{16\beta\kappa^2(t)}{\sigma^2}
                 \exp\left(-\frac{2\kappa^2(t)}{\sigma^2}\right).
\end{equation}
The function
\[
  f(\kappa) := \kappa^2 \exp\left(-\frac{2\kappa^2}{\sigma^2}\right)
\]
has a unique maximum at $|\kappa| = \sigma/\sqrt{2}$, with
\[
  f_{\max} = \frac{\sigma^2}{2e}.
\]
Therefore,
\begin{equation}
  \max_t c(t)
  = c_0 + \gamma \frac{16\beta}{\sigma^2} f_{\max}
  = c_0 + \gamma \frac{16\beta}{\sigma^2} \cdot \frac{\sigma^2}{2e}
  = c_0 + \frac{8\beta\gamma}{e}.
\end{equation}
This shows that the dissipation bump is bounded and independent of the initial condition $\kappa_0$ (provided the trajectory passes through the region where $|\kappa|=\sigma/\sqrt{2}$).

\subsection*{E. Phase-Space Dissipation and Relative Coordinates}

\subsubsection*{E.1 Relative Coordinates}

Define relative coordinates
\[
  q_r = q_1 - q_2, \quad p_r = p_1 - p_2.
\]
The coupled mechanical equations induce a dissipative dynamics in $(q_r,p_r)$.

Without reproducing all algebra, the key qualitative property is:

\paragraph{Proposition E.1 (Dissipative Spiral in Relative Phase Space).}
Under the multiplicity-shaped dissipation $c(t)\ge c_0>0$ and Lyapunov-based coupling dynamics, the relative coordinates $(q_r,p_r)$ exhibit trajectories that spiral toward the origin in the $(q_r,p_r)$-plane.

\emph{Sketch of Proof.} In the symmetric case for masses and appropriate symmetry assumptions on stiffness, the relative mode behaves like a damped oscillator with time-varying damping $c(t)$ bounded away from zero. Standard results for linear systems with positive damping imply asymptotic decay of $(q_r,p_r)$ to $(0,0)$, and trajectories in the phase plane form spirals converging to the origin. Numerical integration confirms this qualitative picture for the benchmark parameter ranges. \qed

\subsection*{F. Summary of Key Bounds and Properties}

For ease of reference, we summarize the principal mathematical properties established in this appendix:

\begin{enumerate}[label=(\arabic*)]
  \item \textbf{Spectral structure:} The contracted mass block $M_{\text{mass}}(\kappa)$ is symmetric with eigenvalues $\lambda_\pm(\kappa)$ and gap $\Delta(\kappa)$ as given above. In the symmetric mass case, $\Delta=2|\kappa|$.
  \item \textbf{Invertibility and conditioning:} If $|\kappa|<\sqrt{m_1 m_2}$, then $M_{\text{mass}}(\kappa)$ is invertible. In the symmetric mass case with $|\kappa|\le \alpha m$ ($0<\alpha<1$), the condition number satisfies $\kappa_2 \le (1+\alpha)/(1-\alpha)$.
  \item \textbf{Lyapunov monotonicity:} The Lyapunov function $H(\kappa)=2\kappa^2$ satisfies $\dot H = -16\beta\kappa^2\le 0$ under the feedback law $\dot\kappa=-4\beta\kappa$.
  \item \textbf{Relaxation time:} The time to reach $H(t)\le \alpha H(0)$ is $t_\alpha \ge \ln(1/\alpha)/(8\beta)$; in particular, $t_{95} = \ln 20/(8\beta)$.
  \item \textbf{Multiplicity dynamics:} The smooth multiplicity observable $M(\Delta)=1+\exp(-\Delta^2/(2\sigma^2))$ has derivative $\dot M = -(16\beta\kappa^2/\sigma^2)\exp(-\Delta^2/(2\sigma^2))\le 0$ and is smooth in $\kappa$.
  \item \textbf{Dissipation bounds:} The dissipation coefficient $c(t)=c_0+\gamma|\dot M(t)|$ satisfies $c(t)\ge c_0$ and, in the symmetric case, is bounded above by $c_0 + 8\beta\gamma/e$.
  \item \textbf{Relative-mode dissipation:} In relative coordinates $(q_r,p_r)$, the system behaves as a damped oscillator with time-varying but bounded damping, yielding spiraling trajectories toward the origin.
\end{enumerate}

These results provide the explicit mathematical guarantees---Lyapunov decrease, operator norm bounds, and dissipation limits---underpinning the numerical and conceptual framework presented in the main article.
\nocite{*}
\bibliographystyle{plain}
\bibliography{references}

\end{document}
